% Generated draft: requires mathematical and source review.
\documentclass[11pt]{article}
\usepackage[margin=25mm]{geometry}
\usepackage{fontspec}
\setmainfont{Latin Modern Roman}
\usepackage{amsmath,amssymb,mathtools,unicode-math}
\setmathfont{Latin Modern Math}
\usepackage{enumitem,needspace,xurl,hyperref}
\setlength{\parindent}{0pt}
\setlength{\parskip}{5pt}
\setlength{\emergencystretch}{3em}
\newcommand{\N}{\mathbb{N}}
\newcommand{\Z}{\mathbb{Z}}
\newcommand{\Q}{\mathbb{Q}}
\newcommand{\R}{\mathbb{R}}
\newcommand{\C}{\mathbb{C}}
\newcommand{\F}{\mathbb{F}}
\newcommand{\A}{\mathbb{A}}
\newcommand{\PP}{\mathbb{P}}
\newcommand{\E}{\mathbb{E}}
\newcommand{\eps}{\varepsilon}
\newcommand{\dd}{\,\mathrm{d}}
\begin{document}
\section*{UnsolvedMath ID: 10400148}
\textbf{Source problem code:} catalogue entry
\subsection{Definitions and mathematical statement}
The source asks to find and classify all modular categories, with completeness and an appropriate criterion for equivalence. Under the customary finite semisimple convention over \(\C\), a modular category is a \(\C\)-linear semisimple rigid ribbon tensor category with simple unit, finite-dimensional morphism spaces and finitely many isomorphism classes of simple objects \(X_1,\ldots,X_s\), whose matrix
\[S_{ij}=\operatorname{tr}_{X_i\otimes X_j}\bigl(c_{X_j,X_i}\circ c_{X_i,X_j}\bigr)\]
is nonsingular. Here \(c\) is the braiding and the trace is the ribbon categorical trace. A classification in this convention describes all such categories up to equivalence preserving the tensor, braiding and ribbon structures. The excerpt supplies no definition, ground field, unitarity condition or chosen equivalence relation, so this familiar definition identifies a conditional interpretation rather than silently restricting the source's ``all.'' No rank or dimension bound is stated.

The extracted text then introduces, for a TQFT \((V,Z)\), the formal generating series
\[P_{(V,Z)}(t)=\sum_{g=0}^{\infty}\dim V(\Sigma_g)\,t^g,\]
where \(\Sigma_g\) is a closed genus-\(g\) surface and the dimensions require finite-dimensional state spaces in the chosen ground field. This is a formal series; no analytic convergence is required. The next sentence says that the \emph{following} problem refines Problem 8.1, but that following problem and Problem 8.1 are absent. Consequently the definition of \(P_{(V,Z)}\) is preserved as adjacent background, without inventing a second classification or a missing condition on its series.
\subsection{Short English statement}
Classify all modular categories up to the intended structure-preserving equivalence. The source does not fix the category conventions; its following TQFT generating-series paragraph is background to an absent next problem.
\subsection{Sources}
\begin{itemize}
\item Local UnsolvedMath registry, permanent ID \(10400148\), exact source code \texttt{catalogue entry}.
\item Ohtsuki, \emph{Problems on invariants of knots and 3-manifolds} (2002), Problem 8.2, PDF page 122: \url{https://msp.org/gtm/2002/04/gtm-2002-04-024s.pdf}. Recorded access 2026-07-29; extraction pdf-font-structure.
\end{itemize}
\end{document}
